\scanpage{159}\mspage{17}
% archive sheet number: 159
Seule la notion a (fbt \inserted{ps.} filtrante à g \deleted{ou à droite}) ne dépend que de $\mathrm{Ord}\,X$ -- mais bien sûr si $X$ est \uuline{préordonnée}, toutes trois ne dépendent que de $\mathrm{Ord}\,X$. Notons que dire que ça ne dépend que de $\mathrm{Ord}\,X$, ou que de $\mathrm{Esp}\unsure{\mathrm{esp}}\,X$, revient au $=$, puisque les deux se déterminent mutuellement. La condition a) signifie que $\mathrm{Esp}\,X$, ou aussi $\mathrm{Top}(X)$ est \uline{irréductible}, \inserted{On dira que $X$ est \uline{irréductible} (\unsure{abrév.} -- \deleted{fbt.} filtr. à g).} les \deleted{autres} autres notions ne \unsure{peuvent} s'exprimer en termes de $\mathrm{Esp}\,X$ \deleted{ou} de $\mathrm{Ord}\,X$ seul, mais elles s'expriment en termes de Topos $X$.

La condition b) ($X$ loc. fbt filtrante) signifie que $\mathrm{Top}(X)$ est \uline{loc irréductible}, et on dira aussi que $X$ est \uline{loc irréductible} (au lieu de loc. fbt. filtrante). Cela signifie que pour tt $X'$ étale sur $X$ et $0$-connexe, $X'$ est fbt filtrant à g. i.e. est irréductible.

Reste à exprimer la condition filtrante elle-même en termes de $\mathrm{Top}(X)$. Je rappelle à ce propos :

\uline{Proposition} 2.3. Les conditions suivantes sur un topos $\unsure{\mathscr X|\mathcal X}$ sont équivalentes (quand $S$ \inserted{(\unsure{comme d'hab.|à un is. près})} désigne une \unsure{partie|petite} génératrice de $\mathrm{Fais}(\mathscr X)=\mathscr X^{\wedge}$\deleted{\unsure{,}}, \uline{formée d'objets $0$-conn.}$^{*}$

% slanted
\margin{(donc $\mathscr X$ est loc. $0$-connexe).}

a) Si $U,V$ sont deux objets $0$-connexes de $\mathrm{Ob}\,\mathscr X$, $U\times V$ est $0$-connexe

b) Si $U$ est un objet $0$-connexe de $\mathscr X^{\wedge}$, alors $U\longrightarrow e$ est ``$W_0$-asph.'', i.e pour \deleted{tout} $V$, la projection $V\times U\longrightarrow\deleted{\unsure{X|U}}V$ induit $\pi_0(U\times V)\xrightarrow{\sim}\pi_0(U)$

% slanted
\margin{Il faudrait expliciter la démonstration, ou celle de 2.9 ci-dessous (p. 25).}

c) Si $U,V$ dans $\mathscr X^{\wedge}$, alors

\[
\pi_0(U\times V)\xrightarrow{\sim}\pi_0(U)\times\pi_0(V)
\]

d) Comme a), avec $U,V$ dans $S$. (Mais \deleted{b) Comme b), avec} \unsure{supposer|supposons} les objets de $S$ $0$-connexes)
